\chapter{全书综合研究工坊}

本工坊不是第十章，而是把前九章的方法合并为一次可以独立完成的探究。读者将从一个现实问题出发，建立证据档案，比较竞争理论，接受反例压力，并提交可审计的结论。项目的成功标准不是找到“绝对正确答案”，而是让任何读者都能看出：问题怎样限定、证据从哪里来、哪一步仍不确定、什么新材料会改变判断。

\section{总案例：一条城市停课预警}

周日晚十点，社交平台流传一张“市教育局通知”截图，称次日因极端天气全市停课。截图带有官方标识和红色印章，多个家长群转发；一个拥有大量关注者的本地账号说“已经向内部人士确认”。天气应用显示暴雨概率为 \(70\%\)，但教育局网站没有通知。生成式搜索助手回答“该市已经宣布停课”，并给出两个新闻链接，其中一个链接只报道邻市停课，另一个页面无法打开。

十一点，教育局值班人员在官方账号发文：“目前尚未作出全市停课决定，请以正式通知为准。”部分家长认为这等于否认停课，另一些人认为“尚未决定”暗示第二天很可能停课。凌晨零点，气象部门把暴雨预警从黄色提升为橙色。校长收到区级工作群消息：“各校做好线上教学准备，是否停课待统一通知。”有班主任为了让家长安排照护，私下说“基本确定不用来”。

凌晨一点，教育局正式宣布部分低洼区域学校停课，其余学校延迟到校。最早截图中的“全市停课”是伪造的，但发布者可能依据一份未定稿内部方案制作；天气风险真实，班主任的预测对本校恰好正确。第二天，有人以“最后确实停了许多学校”为由为截图辩护，也有人把所有早期提醒都称为谣言。

这个案例有意让真假、理由、来源、行动和责任交错。不要立即给人物贴“知道／不知道”的总标签。先完成以下档案。

\section{第一阶段：制作命题与时间档案}

把模糊问题“明天到底停不停课”拆成带量词和时间的命题：

\begin{enumerate}
\item \(p_1\)：周日晚十点时，教育局已经决定全市停课。
\item \(p_2\)：周一全市所有学校事实上停课。
\item \(p_3\)：周一部分学校将停课或延迟到校。
\item \(p_4\)：周日晚十点时，家长有充分理由准备备用照护。
\item \(p_5\)：最早截图是教育局正式发布的文件。
\end{enumerate}

同一证据对这些命题作用不同。天气概率支持风险与准备，不支持文件来源；值班人员的“尚未决定”直接反驳 \(p_1\)，却未反驳 \(p_3\)；班主任可能合理预测本校安排，但无权断言全市决定。若不先拆命题，“最终部分停课”会被错误用来追认伪造截图。

建立逐小时表格，至少包含时间、说者、精确内容、直接来源、转述层级、当时可得反证、行动后果和事后真值。所有评价必须使用当时信息，另设一栏记录事后调查。这样可防止后见偏差：结局为真不等于早先相信方式适当，结局为假也不等于早先行动不合理。

\begin{longtable}{p{0.07\textwidth}p{0.15\textwidth}p{0.24\textwidth}p{0.18\textwidth}p{0.20\textwidth}}
\toprule
时点 & 主体或文档 & 精确命题 & 当时证据状态 & 可合理采取的行动\\\midrule\endhead
22:00 & 最早截图 & 全市已经决定停课 & 来源未认证，只有版式线索 & 不作事实转发，核查官方渠道\\
23:00 & 值班账号 & 尚未作出全市决定 & 官方账号可认证，语义有时间限定 & 保持准备，不把它解释成永不停课\\
00:00 & 气象预警 & 天气风险升级 & 气象来源可靠，但不决定教育政策 & 提高防护，等待学校范围通知\\
01:00 & 正式通知 & 部分停课、其余延迟 & 授权来源与范围明确 & 按所在学校执行并保存更新\\
\bottomrule
\end{longtable}

\section{第二阶段：用九章重新提问}

每一章不是给同一问题换一个术语，而是揭示一组彼此关联、却不能互相替代的失败。按下列次序写九段分析，每段内部要连接本章的几个核心区分，并注明与相邻问题的边界。

\begin{enumerate}
\item \textbf{知识与评价}：分别评价截图发布者、转发家长、班主任和谨慎等待者的真值、理由、过程与责任；再检查班主任的真判断是否具有 Gettier 结构。画出根据路径与真值路径，说明“碰巧说对”为什么不能仅凭真值升级为知识。
\item \textbf{怀疑与运气}：家长是否必须排除截图被完美伪造的所有可能？区分由具体异常触发的怀疑与纯粹逻辑可能，再用邻近情形检验班主任的成功是否安全、是否归功于相关能力或可靠来源。
\item \textbf{理由、语境与击败}：低风险聊天和必须安排照护时，“我知道明天停课”的使用是否变化？把知识归属、断言得体和行动复核分开；判断官方“尚未决定”是反驳型还是削弱型击败，并说明信念资格怎样随时间改变。
\item \textbf{认识来源}：印章、字体和版式提供什么知觉理由，为什么外观相似不能替代来源链？第二天的记忆可能怎样发生来源错配？从以往暴雨停课外推本次需要什么稳定性假设？还应区分对自身焦虑的内省与对政策事实的判断。
\item \textbf{他人、专家与制度}：“内部人士确认”是否可追踪，多个群的转发有多少独立性，气象专家、校长和班主任的专长边界分别是什么？进一步解释信息断层怎样形成，哪些发言者因身份被不当高估或低估，以及制度如何同时容纳准备、保密和纠错。
\item \textbf{概率、模型与决策}：把暴雨概率、停课基础率和新预警作为不同信息，说明为什么 \(70\%\) 降雨不等于 \(70\%\) 停课。若给出数值，列明参考类、似然与依赖假设；再把误报、漏报、等待和照护成本写进决策矩阵。
\item \textbf{知识、理解与智慧的价值}：即使家长不能知道最终政策，理解决策流程、适用范围和更新条件是否仍有价值？比较一条真消息、一个可迁移解释和一项审慎行动各自增加了什么，并说明认识成就怎样分配给个人与信息制度。
\item \textbf{数字环境与人工智能}：生成式搜索为什么会把邻市新闻拼接成本市结论？把链接存在、来源可追溯、内容支持和命题真实分开，追踪排序与摘要怎样压缩不确定性；最后设计一句不会放大未决信息、且给出复核入口的系统回答。
\item \textbf{比较认识论}：选择一个关于“慎言”“正名”或角色责任的古典材料，先作内部解释，再说明它与当代断言规范共享什么问题结构、在哪些概念与历史功能上不可互译；若提出综合主张，明确标为研究者的规范建构。
\end{enumerate}

九段完成后再做综合。若一段分析可以不改一字放入另一段，说明你可能只在重复“证据要可靠”。综合的目标是连接差异：来源认证解决伪造，不解决预测；概率更新说明风险，不决定发布权限；社会制度能分配信息，却仍需个体回应击败者。

\section{第三阶段：建立论证矩阵}

选择一个有限结论，例如：“周日晚十点，家长不应把截图作为事实转发，但可以据天气与多渠道警报准备备用照护。”把它拆为两项行动结论，避免把“不转发”误写成“什么都不做”。

\begin{argumentbox}[结论的基本结构]
\begin{enumerate}
\item 截图的授权来源尚未认证，且官方渠道没有相应决定。
\item 事实性转发会让接收者误以为决定已经作出，错误具有传播外部性。
\item 因此，在完成来源核验前，不应把“全市已停课”作为事实转发。
\item 同时，气象预警、时间压力与照护安排的不可逆成本提供独立的准备理由。
\item 准备备用照护不要求相信全市已经停课，且可在正式通知后调整。
\item 因此，可以采取低成本准备，同时保持对政策命题的判断暂停。
\end{enumerate}
\end{argumentbox}

为每个前提建立矩阵：支持证据、反对证据、来源类型、失败模式、可信程度、若前提为假对结论的影响。至少加入两个最强反对。反对一可说，等待认证会让家长来不及准备；回应应指出准备与事实转发可以分离。反对二可说，快速转发能提醒更多人注意真实天气风险；回应要比较带限定的风险提醒是否实现同样收益而减少误导。

随后进行“最小删除测试”：依次删除每个前提，看结论是否仍成立。若删除一项毫无影响，该项可能是修辞；若一个隐含假设删除后论证崩溃，应把它显式写出。再进行“对手交换测试”：把截图替换为你不信任群体发布的相同截图，看评价标准是否改变。这个测试能发现身份偏见与立场性双重标准。

\section{第四阶段：设计区分竞争理论的材料}

好的研究不只用案例说明自己喜欢的理论，还要让竞争理论给出不同预测。为总案例制作下列最小变体：

\begin{enumerate}
\item \textbf{只改来源}：文本完全相同，但一个来自可验证官方签名，一个来自匿名截图。
\item \textbf{只改真值}：发布过程相同，一版最终全市停课，一版天气转向而不停课。
\item \textbf{只改利害}：一位家长无需安排照护，另一位若午夜前不请假就会失去工作。
\item \textbf{只改注意}：一位读者看见“草案”水印，另一位因图片裁切看不见。
\item \textbf{只改独立性}：五个账号独立联系不同学校，或五个账号都转载同一人。
\item \textbf{只改算法介入}：搜索系统只是列出链接，或直接生成“已经停课”的摘要。
\end{enumerate}

对每组分别记录知识判断、置信度、可断言性、可行动性与责备程度。若风险只改变行动而不改变知识，支持知识与行动分离；若来源改变知识却不改变结局，说明真值不足；若五个独立来源显著提高支持而转载不提高，证言网络的独立性得到体现。数据不能机械决定理论，但会迫使理论解释完整模式。

若邀请他人判断，先写清参与者看到的材料、问题顺序和术语含义。不要把“同意程度”直接当成概念真值。开放回答可揭示参与者究竟根据来源、后果还是道德责备作答。报告不只给平均数，也列出主要解释类型与异常反应。

\section{第五阶段：证据与引文审计}

每个事实性断言标一种证据角色：

\begin{description}
\item[原始记录] 官方通知、原图、时间戳、版本历史和完整对话。
\item[当事证言] 发布者、值班人员、家长与校方的陈述，注明取得时间。
\item[方法资料] 平台传播机制、来源认证规范或风险管理文件。
\item[哲学原典] 支持理论重构的著作、论文和古典文本。
\item[二手研究] 帮助解释争论史、译法或经验结果的学术材料。
\end{description}

原始记录能证明某条消息何时出现，不自动证明消息内容；当事人回忆能说明其经历，也可能受事后信息影响；方法标准说明应怎样处理，不证明本案例已经遵守。每条引文旁写一句“它在这里证明什么”，若写不出，就可能只是权威装饰。

直接引语必须逐字核对并给定位信息，省略要标明且不改变语义；转述同样需要来源，但应用自己的句法和论证结构。古典材料需标篇章或章句与所用译本，现代论文需核对作者、题名、年份、卷期、页码或 DOI。网页材料记录发布机构、标题、版本、日期和访问日期。无法确认版本时降低主张强度。

做一次“反向引文审计”：不是从正文找参考文献，而是从参考文献逐条回正文，检查是否真的支持相邻主张。标题相关不等于内容支持；引用一篇批评文章时，不能把它批评的观点归给作者本人。若二手文献引用原典，应尽量回到原典，同时保留二手研究对解释的贡献。

\section{第六阶段：从结论到制度建议}

哲学分析若转化为建议，必须说明建议对应哪种故障。针对伪造，使用可验证官方渠道和来源记录；针对范围歧义，把“全市”“部分”“待定”作为强制字段；针对转述失真，让平台显示首发来源与修改；针对时间压力，允许学校发布分阶段状态；针对照护不平等，提供不依赖提前内部消息的支持措施。

不要以“提高媒体素养”承担所有修复。个体查验有成本，官方延迟、平台界面和组织保密规则会塑造错误。也不要以制度问题免除个体责任：明知来源不明仍删去限定转发，与谨慎提示尚未认证不同。建议应分别列个人、学校、教育局、平台和工具提供者可以控制的环节。

每项建议附四个验收问题：

\begin{enumerate}
\item 它减少哪一种具体错误，而不是笼统“增强信任”？
\item 它是否产生新成本，例如隐私、延迟、警报疲劳或权力集中？
\item 谁负责实施，谁有权暂停，谁能够申诉？
\item 用什么模拟或记录证明它在下一次事件中有效？
\end{enumerate}

建议的认识质量不由善意保证。来源徽章可能被用户误读为内容真实，过多正式更新可能淹没关键变化，强制实名可能让弱势爆料者沉默。每项修复都要接受反例与分布性影响检查。

\section{三类可替换的结课课题}

若不使用停课案例，可从以下课题任选一项。每项都必须保留时间档案、九问、论证矩阵、竞争理论变体与证据审计。

\subsection{医疗 AI 与第二意见}

系统把影像标为高风险，首位医生同意，第二位医生反对；病理结果尚未取得。研究可问：患者现在知道什么？模型概率如何与行动阈值关联？医生是否独立于模型？不同群体上的误差怎样改变信任？知情同意需要说明哪些不确定性？比较材料可围绕专家角色、审慎与行动责任，但不得用古典格言替代临床证据。

至少制作三个数据情景：模型与医生错误独立、医生先看到模型而被锚定、模型部署人群偏离验证样本。相同“二对一”意见在三种情景中证据权重不同。制度建议要具体到读片顺序、覆写记录、患者申诉和停用阈值。

\subsection{争议历史记忆与公共纪念}

社区对一次旧事件有冲突记忆：档案不完整，口述史彼此不同，官方纪念文本只保留一方。研究可问：记忆鲜明度与来源资格如何区分？沉默是没有证据还是证言被窒息？群体能否拥有无人完整掌握的历史知识？公共文本应如何报告不确定性？

材料至少包含一份同期记录、一份晚年回忆、一项制度档案缺口和一项后来的政治使用。不得以档案缺失自动否定口述，也不得以受害经验自动决定所有时间细节。建议需同时处理事实核验、第一人称尊重和诠释参与。

\subsection{科学争议与公共传播}

选择一项仍有模型不确定性的科学政策问题。研究可问：专家共识如何形成？公众怎样识别元专长？概率区间怎样转化为决定？媒体何时把分歧夸大成“科学家各说各话”，又何时把真实少数意见压平？数字推荐怎样改变证据可见性？

至少区分事实分歧、模型分歧与价值阈值分歧。建立来源谱系，检查多个报道是否来自同一新闻稿。制度建议不应要求科学家等到完全确定才发言，而要设计分级断言、更新说明和更正可见性。

\section{最终交付包}

完整项目包含七份材料：

\begin{enumerate}
\item 一页问题与范围声明，列研究问题、核心命题和不处理的相邻问题。
\item 一份不少于十个节点的时间—来源档案。
\item 一张论证图和一张竞争理论矩阵。
\item 六个最小案例变体及判断记录。
\item 一份带证据角色说明的书目与引文审计表。
\item 一篇八千至一万二千字论文，正文明确区分事实、解释与规范建议。
\item 一页修订报告，列出最强反对、实际改动与仍未解决的问题。
\end{enumerate}

论文摘要应能回答五句：我研究什么有限问题；现有两种最佳答案是什么；我用什么案例或证据区分它们；我的结论是什么；结论在哪个条件下会失效。若摘要只能说“本文讨论知识的重要性”，说明范围仍过宽。

\section{口头答辩与失败测试}

十分钟答辩不重复论文目录。前三分钟重构对手的最强论证，接着三分钟说明自己的区分性证据，再用两分钟陈述结论和边界，最后两分钟主动提出最危险反例。答辩者若只防守，不显示什么会改变自己，不能证明认识自主。

答辩同伴从四类问题中各选一项：

\begin{enumerate}
\item \textbf{概念问题}：你的“知道”“合理”或“理解”在全文是否保持同义？
\item \textbf{证据问题}：哪条关键证据若被撤回，结论变化最大？
\item \textbf{对手问题}：竞争理论能否吸收你的案例，只付出很小代价？
\item \textbf{责任问题}：你的建议把核验成本交给了谁，是否公平且可执行？
\end{enumerate}

答辩后必须修改至少一处结构，而不只是润色。可以缩小结论、补上桥梁前提、承认两种评价分离、替换不支持的来源或删除过度类比。修订报告要引用改动前后句子并说明理由。

最后进行“彻底失败测试”。假设主要结论为假，写出最可能的世界：是新证据证明来源真实，还是行动代价使不同策略更合理，抑或比较桥梁根本不成立？然后问现有研究为何尚未排除该世界。能够清楚描述失败条件，不削弱论文；它表明结论的强度已经与证据相称。

\section{完成标准}

项目完成不等于材料齐全。以下条件必须同时满足：

\begin{itemize}
\item 没有用事后真值替代当时认识资格，没有用主观无责替代知识。
\item 每个重要来源可定位，直接引语已核对，链接与相邻主张真正对应。
\item 至少一个案例只改变一个变量，并确实让竞争理论给出不同解释。
\item 比较部分标明内部解释、结构比较与作者建构，不把相似当同一。
\item 概率数字带分母、参考类和依赖假设，行动建议另行说明价值权衡。
\item 社会与数字分析追踪信息链，不以“个人要谨慎”结束制度问题。
\item 结论明确列出适用范围、残余不确定性和未来修正触发条件。
\end{itemize}

达到这些标准时，读者获得的不只是一篇课程论文，而是一套可迁移的认识实践：把大问题拆成命题，把证据放回时间与来源，把理论置于反例压力，把依赖转化为可问责关系，并在不得不行动时仍保留修正的通道。这正是一本问题驱动认识论教材最终要培养的能力。

\index{研究设计}
\index{证据审计}
\index{来源谱系}
\index{最小对比}
\index{认识责任}
\index{问题驱动}
